% ===========================================================================
% TMLR submission, paper 2 of a two-paper split (round 22).
% Preamble and class setup copied from the companion paper's main.tex.
% Sections follow the body budget of ROUND22_PLAN.md, "Paper 2 spec".
% Every number is carried over from the companion paper's text or from the
% frozen arm reports in the corpus; none is recomputed here (NUMBERS.md).
%
% Compile: pdflatex -> bibtex -> pdflatex x2, then the same four for
% supplement.tex (the supplement reads main.aux through xr-hyper).
%
% Camera-ready: switch to \usepackage[accepted]{tmlr}, restore the real
% author block, set \month/\year/\openreview below.
% ===========================================================================
\documentclass[10pt]{article}
\usepackage{tmlr}
% If accepted: \usepackage[accepted]{tmlr}
% Preprint (arXiv): \usepackage[preprint]{tmlr}

\usepackage{hyperref}
\usepackage{url}
\usepackage{amsmath}
\usepackage{amssymb}
\usepackage{booktabs}
\usepackage{array}
\usepackage{longtable}
\usepackage{microtype}
\usepackage{graphicx}
\raggedbottom

% zero-width, hyphen-free break opportunity, used inside long hex digests
\newcommand{\bk}{\discretionary{}{}{}}

\title{Characterize Before You Loop: What Transfers From a Zero-Shot
Proposer's Template and What Does Not}

% DOUBLE-BLIND: real author block intentionally absent from this file.
\author{\name Anonymous authors \email anonymous@example.org \\
      \addr Withheld for double-blind review}

\def\month{MM}  % camera-ready
\def\year{YYYY} % camera-ready
\def\openreview{\url{https://openreview.net/forum?id=XXXX}} % camera-ready

\begin{document}

\maketitle

\begin{abstract}
A weak-tier language model proposer, asked in a zero-shot call to pack circles in the unit
square, concentrates on a grid-plus-filler template. A companion submission gives the closed
form naming its modal output. That template is the anchor, where its outputs settle at a
given size. Its recipe family has a largest member there, the family argmax. This
paper measures which survives a change of setting.
Eight registered arms change four things. They vary the parent, the
iteration, the serving path with its reasoning setting, and the vendor. Under a better
in-family parent, on the agent-runtime instrument, 26 of 26 valid outputs keep or improve that
parent's score, and none returns to the anchor. Under its own anchor as parent the model leaves that anchor,
missing the registered bar. Iteration erodes it across five archive-and-mutate
generations on that instrument. The same prompts on a pinned
path return 0 of 105 valid packings. Extended thinking restores validity at the one cell
tested. At a second vendor the registered pooled bar is met; a post hoc split shows the cells
that test anchoring miss it. Under an instruction to improve, the anchor does not persist. The family argmax does, as
a ceiling no arm clears: no output valid at $10^{-9}$ exceeds it by more than $10^{-6}$.
Clearing it would mean leaving the recipe family,
and no valid output did. In plain terms, the anchor value is easy to dislodge and the
ceiling above it is not. Across every prompt, serving path and vendor tried, no valid packing
cleared the best member of that recipe family: that member is the ceiling, and the ceiling
transfers.
\end{abstract}

\section{Introduction}\label{sec:intro}

Ask a language model to place 21 circles in a unit square to maximize the sum of their radii.
It will not search. It reaches for the nearest square grid, five by five ($k = 5$) at
radius $1/(2k)$, and truncates it, leaving four cells empty. A companion submission
\citep{companion} gives that behaviour a closed form, naming the modal output in advance.
It reports that form as the empirical mode at all seven $N$ tested, at four of the five
held out, and under a perturbed container and paraphrase. The two cell sets are listed in
Table~\ref{tab:glossary}. This paper's own generation-0 rows were sampled
fresh under its loop arm, on the same instrument and prompt bytes
\citep{companion}. At $N = 31$ they put the anchor on 7 of 8 valid
outputs. At $N = 13$, on 4 of 9 (Table~\ref{tab:overview}).

A characterization that sharp invites an inference. What a proposer emits zero-shot is what it
will contribute inside a discovery loop, so the loop need not run. Discovery systems in the
FunSearch line \citep{romeraparedes2024funsearch,novikov2025alphaevolve} do not call their
proposer zero-shot. They condition it on a parent and iterate it against an archive. Each is a
change of setting, as is the serving path \citep{companion}. Each can break the
characterization.

\textbf{Objects and arms.} A \emph{template} is one
member of the grid-plus-filler recipe family. The \emph{anchor} is the template the model's
outputs concentrate on at a given $N$. The \emph{family argmax} is the largest value the
family can express at that $N$. A \emph{discriminating cell} is an $N$ where the two differ,
so only there can anchoring be told from family search. Table~\ref{tab:glossary} carries the
rest. Eight registered arms change four things. Arm MU changes
the parent. Arm L adds iteration, feeding a scored proposal back as the next parent. Arms P
and P-D change the serving path and the reasoning setting on it. Arms GM, GM2, GM3 and V
change the vendor, GM2 and GM3 differing only by output budget.

\textbf{Claim.} A weak-tier LLM proposer's zero-shot anchor in circle packing is contingent on
the serving path and its request parameters. A zero-shot
characterization therefore does not establish that its anchor transfers into a discovery
loop, absent a measurement there. The anchor does not transfer. The family argmax does, as a
ceiling. Under the clearance rule of \S\ref{sec:method}, no output valid at $10^{-9}$
exceeds it by more than $10^{-6}$ in any arm. At the second vendor's discriminating cells
most valid outputs sit on it. Two kinds of residual are reported and neither is a clearance.
The primary $10^{-6}$ tolerance admits circle overlaps up to
that size, and under it 18 arm MU outputs sit above the argmax by at most
$1.53\times10^{-6}$; all 18 fail at $10^{-9}$
(\S\ref{sec:arm-mu}). In arm GM3, 29 outputs valid at $10^{-9}$ sit above the argmax by at
most $5.3\times10^{-10}$, the rounding of a printed coordinate (\S\ref{sec:gm3-v}). Two arms carry no row in Table~\ref{tab:overview}. Arm
GM finished UNDERPOWERED whether its routed rows are counted or not. Arm GM2 returned nothing
parseable at the smaller output budget. Arm GM3 diagnosed the cause at the larger.

\begin{table}[!ht]
\caption{Six of the eight arms: what each changed, its bar, the outcome
and the verdict. Codes like P-MU1 name a registered prediction,
written out in its arm's section and Supplement~S1. L1 to L4 are the loop arm's
four, named in \S\ref{sec:arm-l}. Arms MU and L ran on the agent-runtime instrument, P and
P-D on the pinned path. Bracketed ranges are 95\% Wilson intervals on the rate they follow, reported as descriptive precision; no verdict turns on one.}
\label{tab:overview}
\begin{center}
\footnotesize
\setlength{\tabcolsep}{3pt}
\begin{tabular}{@{}>{\raggedright\arraybackslash}p{0.05\linewidth}>{\raggedright\arraybackslash}p{0.15\linewidth}>{\raggedright\arraybackslash}p{0.24\linewidth}>{\raggedright\arraybackslash}p{0.22\linewidth}>{\raggedright\arraybackslash}p{0.21\linewidth}@{}}
\toprule
Arm & What changed & Registered bar & Outcome & Verdict \\
\midrule
MU & The parent the proposer is conditioned on & Under a better in-family parent, on-prediction rate
at most 20\% and keep-or-improve at least 50\% & 26 of 26 keep or improve [87--100\%], 25 within $10^{-6}$ of the parent's score and one
$1.2\times10^{-6}$ above it, so the on-prediction rate there is 0 of 26 [0--13\%]; under the model's
own anchor, 7 of 28 keep it [13--43\%] and 17 move above it & The rule reduces to one
condition under the better parent; the own-anchor rate misses P-MU1's 50\% bar, so the
anchor is not self-stable; the falsifier is that rule's negation and adds no second result \\
\addlinespace
L & Five archive-and-mutate generations & Pooled conditioned on-prediction rate below the
generation-0 rate in both cells (L2); the falsifier is its negation & $N = 13$: 4 of 9 at generation 0, 5 of
20 conditioned. $N = 31$: 7 of 8, then 6 of 29. Pooled 11 of 17 [41--83\%], then 11 of 49
[13--36\%] & L2 holds; L1 and L3 not met; L4 open (0 of 49 clear the family, [0--7\%]) \\
\addlinespace
P & The serving path: the bare pinned chat-completions endpoint & A floor of five valid
outputs per cell & 0 of 105 valid at the square cells [0--4\%]; 4 of 75 at the held-out cells
[2--13\%] & Every direct-emission prediction UNSCOREABLE: no cell reaches the floor \\
\addlinespace
P-D & Extended thinking on against off, same pinned path & The reasoning budget is the cause,
against the missing context, against anything else & Thinking on, 12 of 14
valid with the anchor the plurality value at 5 of 12 [19--68\%]; off, 0 of 13 [0--23\%] &
The reasoning-budget hypothesis holds; extended thinking is the component \\
\addlinespace
GM3 & A second vendor's gemma at a 16,384-token budget & Pooled rate at least 30\%, and the
anchor modal in 5 of 5 scoreable cells & 46 of 80 pooled (57.5\%, [47--68\%]); modal
in 2 of 5; 11 of 43 at the discriminating cells (26\%, [15--40\%]) & The pooled bar is met;
the discriminating cells, the only ones that test anchoring, miss it (post hoc split) \\
\addlinespace
V & Thirteen free-tier aliases & At least 10 of a model's 25 invocations valid, then V1 at
least 50\% and V2 at most 10\% & Ten aliases below the floor; two aliases meet the V1 pooled
bar at point estimate, on 8 of 12 [39--86\%] and 11 of 18 [39--80\%] of their valid outputs,
and a third does not, at 0 of 15 [0--20\%] & The boundary is located, not attributed \\
\bottomrule
\end{tabular}
\end{center}
\end{table}

\textbf{Contributions.}
\begin{enumerate}
\item \emph{A measurement of transfer, not an argument about it.} Each change of setting got
its own arm, with thresholds and a falsifier fixed before the first
invocation. Every deviation is tabled.
\item \emph{A named instrument component, at one cell.} On the pinned path at $N = 13$, arm
P's validity collapse survives the one system line tested, a request for care. One
registered condition (P-D, whose rows both submissions report) and two post hoc diagnostics
separate the candidates. Extended thinking carries the anchor.
\item \emph{A cross-vendor boundary that is located but not attributed.} Two second-vendor
measurements disagree, with model variant and serving path confounded.
\end{enumerate}

\textbf{Scope.} The proposer is one vendor's weak-tier alias
(\texttt{claude-haiku-4.5}) except where another is named. The conditioning arm runs
at three trap cells, $N = 13$, 21 and 31. The loop arm runs at two of them.

\emph{What is not claimed.} Nothing about the companion paper's numerical baseline, about
which constructions beat the family argmax, or about its
mid-tier alias. Those are its results, cited and not restated \citep{companion}.
Nothing about mechanism inside the weights, about what a discovery system does at
generation zero, or whether diversity machinery
\citep{mouret2015mapelites,lehman2011novelty} repairs or explores. Preregistration
is an adopted standard, not a contribution (\S\ref{sec:related}).

\section{The template and its closed form}\label{sec:task}

This section restates the object under test. Every definition below is the companion paper's
\citep{companion}. Only the re-score at the end is this paper's own.

\textbf{The benchmark.} Maximize the sum of radii of $N$ non-overlapping circles in the unit
square. Feasibility is decidable exactly from emitted coordinates, so scoring needs no human.
The direct-emission prompt fixes the count, states containment and non-overlap, forbids
code, and demands a raw Python list of \texttt{[x, y, r]} triples. The program channels
replace that clause with a program contract. The prompt and its digest are in the
supplement.

\textbf{The recipe family.} The \textbf{recipe family} is the parametric set of
grid-plus-filler constructions. A \textbf{template} is one member of it. The \textbf{anchor}
is the template the distribution concentrates on at a given $N$.

A $k \times k$ grid places one circle per cell center with $r_{\mathrm{grid}} = 1/(2k)$. Then
$m$ fillers sit on interior grid vertices, tangent to the four grid circles around them,
with $r_{\mathrm{filler}} = (\sqrt{2} - 1)/(2k)$ and $0 \le m \le (k-1)^2$. Hence
$V(k, m) = k/2 + m(\sqrt{2} - 1)/(2k)$.

When $N < k^2$ a smaller complete grid is available, the $(k-1) \times (k-1)$ lattice at
radius $1/(2(k-1))$ with fillers on its interior vertices. The model does not fall back to
it. It \emph{truncates} instead, laying out the $k \times k$ lattice, occupying $N$ cells and
keeping the radius, which gives $T(k, N) = N/(2k)$.

\textbf{The order rule and the branch rule.} Write $k^*(N) = \mathrm{round}(\sqrt{N})$ for the
nearest-square order. This is a definition, not a hypothesis. The falsifiable content is the
\textbf{branch rule}. Where $k^{*2} \le N$ the model extends, adding $m = N - k^{*2}$ fillers:
the converge branch. Where $k^{*2} > N$ it truncates the $k^*$-grid: the trap branch.
Substituting $k^*$ into the branch condition gives the \emph{trap zones}
$N \in [k^2-k+1, k^2-1]$, which over the sweep range are $[13,15]$, $[21,24]$, $[31,35]$,
$[43,48]$ and $[57,63]$.

The \textbf{family argmax} at $N$ is the largest value the family can express there,
recomputed in closed form over every admissible $(k, m)$. The scripts enumerate every order
$k$ with an admissible filler count, $N - k^2 \le (k-1)^2$, and every truncated order. The
argmax is a value; a packing
that reaches it is a construction. At the trap cells studied here that enumeration returns
the larger
of two candidates: the drop-and-fill $V(k^*{-}1, N - (k^*{-}1)^2)$ and the truncated anchor
$T(k^*, N) = N/(2k^*)$. At $N = 35$ the truncated anchor is larger, so 35 is not a
discriminating cell, an $N$ where the anchor and the family argmax differ. At a converge cell the extended anchor is the argmax, so every
discriminating cell is a trap cell. The anchor there is the truncate
branch $T(k^*, N) = N/(2k^*)$.

\textbf{Independent re-score.} A second implementation recomputed every cell-level count in
this paper, the validity, on-prediction, clearance and row counts per cell
(\texttt{arm\_transfer\_\bk{}independent\_\bk{}rescore.py} in the artifact). The authors wrote
it for this submission, with its own parser, its own validity scorer and a linear program
blind to the recipe. It imports nothing from the arms' scoring code. It re-derived the
closed-form value and the family argmax at all twelve $N$, agreeing to at most
$3.6 \times 10^{-15}$. Of 135 cell-level counts across all arms, 129 agree. The six that
differ are three arm L clearance counts at $N = 31$, two arm GM validity counts and arm P-D's
sampled-row count, each with its cause in Supplement~S1. The re-score also found an arm MU
bucketing error, logged as item 36 of the artifact's corrections ledger, and
\S\ref{sec:arm-mu} carries the corrected figures.

\textbf{Notation and conventions.} Table~\ref{tab:glossary} restates the companion paper's
terms.

\begin{footnotesize}
\setlength{\tabcolsep}{3.5pt}
\begin{longtable}{@{}p{0.26\linewidth}p{0.68\linewidth}@{}}
\caption{Terms used throughout this paper, each used in this sense only.}
\label{tab:glossary} \\
\toprule
Term & Meaning \\
\midrule
\endfirsthead
\toprule
Term & Meaning \\
\midrule
\endhead
on-prediction & emitted sum within $2\times10^{-3}$ of the anchor \\
family argmax & the highest-valued member of the recipe family at that $N$: at the trap cells studied here, the drop-and-fill $V(k^*{-}1, m)$ wherever it beats the truncated anchor. This paper gives it no other name \\
square, held-out and code cell & the three cell sets: a \emph{square cell} is an $N$ addressed with the bare direct-emission prompt, $N = 13$, 17, 21, 31, 35, 37, 43; a \emph{held-out cell} is an $N$ the companion paper kept out of its fit and this paper reuses, $N = 50$, 58, 62, 65, 75; a \emph{code cell} is an $N$ addressed under the program contract instead, $N = 13$, 21, 31 \\
discriminating cell & $N$ where anchor $\ne$ family argmax; only these can distinguish anchoring from family search \\
validity & non-overlap and containment at $10^{-6}$ (primary) and $10^{-9}$ (logged); the $10^{-9}$ count is reported wherever a clearance is read \\
trap cell & an $N$ inside a trap zone, where the branch rule truncates; the discriminating cells are the trap cells studied here whose truncated anchor sits below the family argmax \\
tier & nominal capability class of the serving alias addressed: here the weak tier is the vendor's Haiku alias (\texttt{claude-haiku-4.5}) \\
instrument & the proposer's serving path with the request parameters it carries, the reasoning setting included. Two are used here, an agent runtime and a pinned API, described in \S\ref{sec:method}; they are never pooled, and the word is used in this sense only \\
effort & the size of the reasoning budget a request asks the vendor for. Where a request names none the vendor's routing layer supplies its own default, which is the setting arm P-D switches on \\
quota & the vendor-side cap on requests per interval; a key that reaches it returns no completions until the interval resets \\
UNSCOREABLE, UNDERPOWERED, BELOW-FLOOR & a cell below its arm's registered floor of valid samples; a second-vendor chain arm with fewer than 5 scoreable cells; an arm V alias with fewer than 10 of 25 registered invocations valid. Each claims nothing \\
evaluability floor & the minimum valid-sample count a cell needs before its verdict counts, fixed at each arm's own registration and \emph{not} common across arms; the per-arm values are in \S\ref{sec:method} \\
anchor & the template the proposer concentrates on at $N$, with its closed-form value. That value is the extended anchor $V(k^*, N - k^{*2})$ at a converge cell and the truncated anchor $T(k^*, N) = N/(2k^*)$ at a trap cell. At a discriminating cell it sits below the family argmax \\
$k$-inheritance & an output whose dominant circle radius gives the same grid order $k$ as its parent's; counted per parent condition in the arm MU table \\
located but not attributed & a difference measured between conditions differing in more than one way, so that the difference is established and its cause is not \\
\bottomrule
\end{longtable}
\end{footnotesize}

\section{Registration, instruments, scoring, floors}\label{sec:method}

\textbf{Registration protocol.} For each cell the prompt is pinned and its SHA-256 digest
written to disk first. The $N = 13$ square prompt has a digest beginning \texttt{32db485b},
given in full in the supplement. Predictions live in the header of each arm's analysis script or in
a standalone registration file. Competing predictions, numeric thresholds, tie conventions and
a falsifier are fixed before the first invocation. Raw outputs are stored verbatim, failures
included. Scoring is deterministic and local. Each arm writes its rows to a JSONL ledger.
The supplement carries every arm's ledger, every departure from a registration, and the
per-cell tables for the serving-path and second-vendor arms. Every reported verdict was
registered before its arm's first scored run. Post hoc readings are labelled where they
appear and listed together in \S\ref{sec:limits}.

\textbf{Two instruments, never pooled.} The first is an agent runtime. It does not expose
decoding parameters, so every runtime row is a distributional claim over an observed sampling
regime. It dispatches each call to a \emph{subagent}, its worker process. That process
inherits instruction files outside the task prompt. The digests do not lock those files, and
a replicator cannot reconstruct them. The runtime also publishes no weight hash, so the
alias-to-weights binding is a vendor promise. The second is a pinned API, the vendor's
chat-completions endpoint reached through OpenRouter, at temperature 1.0 with no system
prompt. It exposes and pins temperature and carries no inherited context. Arms MU and L run
on the runtime. Arms P and P-D run on the pinned path. Arms GM, GM2 and GM3 use a second
vendor's first-party API, and arm V uses free-tier aliases at seven attributed vendors. No
count in this paper pools across instruments.

\textbf{Scoring conventions, fixed in advance.} Validity is scored at $10^{-9}$ and
$10^{-6}$; $10^{-6}$ is primary, and the $10^{-9}$ count appears wherever a clearance is
read. Proposers print six to eight decimals, so an
eight-decimal tangency misses contact by about $5\times10^{-9}$. The primary threshold sits far
below the gap of about $10^{-2}$ between the competing constructions a cell can express. Value
matching uses a $2\times10^{-3}$ window. That window never carries a clearance verdict. One
clearance rule serves every arm: an output clears the family argmax when it is valid at
$10^{-9}$ and its verified sum exceeds the closed-form argmax by more than $10^{-6}$. An output
valid only at $10^{-6}$ cannot clear, whatever its sum, and a residual under $10^{-6}$ is
rounding, not clearance. Both kinds are counted and reported where they occur. That argmax is
the family's maximum by construction,
so a clearance is an output outside the recipe family. Completions are parsed with \texttt{ast.literal\_eval}
after stripping code fences. One that fails to yield a list of $N$ strictly positive numeric
triples is a logged parse failure, scored invalid. Every on-prediction rate is computed over
valid outputs and so conditions on execution success. Per-cell validity denominators are
reported beside every rate, except in arm MU. That arm registered three rates, the
on-prediction rate (its registration calls it the anchor-rate), keep-or-improve and $k$-inheritance, pooled across its three cells. Only
P-MU4 (\S\ref{sec:arm-mu}) was registered per cell. Its table therefore gives
pooled counts for those three and per-cell counts for P-MU4.

\textbf{Evaluability floors, per arm.} The evaluability floor is the minimum
valid-sample count a cell needs before its verdict counts, fixed at each arm's own
registration. The second-vendor chain uses 3: a cell below it is UNSCOREABLE for mode analysis
and claims nothing in either direction. A chain arm with fewer than 5 scoreable cells is
UNDERPOWERED. Arms P and P-D use 5. Arm V uses a per-model floor of
at least 10 of a model's 25 registered invocations valid at $10^{-6}$. A model below it is
labelled BELOW-FLOOR and carries no anchoring claim in either direction. Arm L's floor is a
two-part rule on each lineage, the chain of generations grown from one cell
(\S\ref{sec:arm-l}). A lineage needs at least 3 valid proposals at generation 0, the
unconditioned start, and at least 1 in at least 3 of the 5 generations conditioned on a
parent. All four lineages clear it. The weakest, lineage 13-GREEDY at $N = 13$, clears on 5
valid at generation 0 and on valid outputs in 3 of its 5 conditioned generations. Its 4 valid
conditioned proposals are not the quantity the floor tests.

\textbf{The eight arms.} All eight were run once, in one registered corpus. Arms P and P-D
answer an instrument question in the companion and a transfer question here, so both papers
report the same rows in full \citep{companion}. The other six are reported here and appear in
the companion only as a pointer.
Each arm's design is stated where the arm is reported. Supplement~S1 tables all eight.

\section{One conditioning step: arm MU}\label{sec:arm-mu}

Arm MU hands the proposer one parent and one instruction, with no iteration. Given a
parent at the family argmax the proposer keeps it, 26 of 26, and none of those outputs sits
on the anchor. Given its own anchor as parent it leaves, 17 of 28 moving above it.

\textbf{Design.} Each call is the bare prompt, an existing packing, and the instruction to
increase the sum of radii.
The \emph{anchor parent} is built at the model's own unconditioned anchor value $T(k^*, N)$.
The \emph{argmax parent} is built at the family argmax $V(k^*{-}1, m)$, which at these cells
is provably better than the anchor. The \emph{off-family parent} is a diagonal-chain packing
scoring far below either. The registration and the supplement name them A\_anchor, B\_rival
and C\_offfamily. The cells are $N = 13$, 21 and 31. Three parent conditions in three cells,
at $n = 15$ samples each, give 135 invocations.

\textbf{Registered predictions.} P-MU1 predicted an on-prediction rate under the anchor parent
of at least 50\%. P-MU2 predicted an on-prediction rate under the argmax parent of at least
35\%. P-MU3 predicted $k$-inheritance under the argmax parent below 50\%. P-MU4 predicted that
under the off-family parent the modal valid output sits on the $k^*$-grid, scored per cell.
The registered decision rule calls anchoring dissolved when the pooled on-prediction rate
under the argmax parent is at most 20\% and the keep-or-improve rate is at least 50\%. The
falsifier fires on the same condition.

\textbf{The rule's two conditions under the argmax parent.} Both are met, and they are not
independent. The pooled on-prediction rate is 0 of 26 valid outputs. That meets
the rule's 20\% ceiling and fails P-MU2's 35\% floor. Keep-or-improve is 26 of 26, above
its floor of 50\%. The argmax parent outscores the anchor
$T(k^*, N)$ by 0.151 at $N = 13$. The gap is 0.159 at $N = 21$. It is 0.165 at $N = 31$. Each
gap is more than seventy times the $2\times10^{-3}$ matching window, so an output that keeps or
improves on that parent cannot sit on the anchor. The 26 of 26 entails the 0 of 26,
and the rule reduces to one condition.

Keep-or-improve is almost entirely keep. Of the 26, 25 return the parent's score within
$10^{-6}$, and one exceeds it by $1.2\times10^{-6}$. In
coordinates, 18 reproduce the parent to within $10^{-6}$ on every coordinate. The
$1.2\times10^{-6}$ excess is among them, from radii perturbed below that tolerance. One
further output is a mirror image of the parent. The other 7 relocate
filler circles while keeping the radius multiset and the parent's sum to seven decimals. All
26 inherit the parent's grid order too, so P-MU3 is disconfirmed.

The independent evidence is the anchor parent. Its on-prediction rate is 7 of 28,
below P-MU1's 50\% threshold, so the anchor is not self-stable as its own parent. Here
the model departs: 17 of 28 move above the anchor and 4 fall below it. Grid order is kept by
22 of 28, 79\% (60--90\%). The 6 that change it all move above the anchor, and they include
3 that reach the family argmax value. The falsifier negates the persistence prediction
over this pair of conditions, and no other was registered.

\textbf{The unconditioned rate on the same instrument.} Arm L's generation 0
(\S\ref{sec:arm-l}) supplies the no-parent reference on the same prompt bytes, channel and
wrapper. It puts 4 of 9 valid outputs on the anchor at $N = 13$. It puts 7 of 8
there at $N = 31$. The anchor parent gives 4 of 12 and 2 of 11 at those cells. Arm
L did not run at $N = 21$, so no unconditioned row exists there.

\textbf{P-MU4, held in two cells of three.} The anchor's grid order does reassert itself in the one
condition registration predicted. Given the off-family diagonal parent, 20 of 34
valid samples abandon the parent and emit a $k^*$-grid construction. The prediction holds
at $N = 21$ and at $N = 31$, and not
at $N = 13$.

\begin{center}
\begin{small}
\setlength{\tabcolsep}{3.5pt}
\begin{tabular}{@{}p{0.18\linewidth}cp{0.26\linewidth}p{0.15\linewidth}p{0.21\linewidth}p{0.06\linewidth}@{}}
\toprule
Parent & Valid & On-prediction (95\% Wilson) & Registered threshold & Score against parent & $k$-inher. \\
\midrule
None (arm L generation 0, same instrument) & 17 & 4/9 at $N = 13$ (44\%, [19--73\%]); 7/8 at $N = 31$ (88\%, [53--98\%]); no row at $N = 21$ & unconditioned reference, no bar & --- & --- \\
Argmax parent (a packing at the family argmax) & 26 & 0/26 (0\%, [0--13\%]) & dissolved if $\le$ 20\% & 26/26 (100\%, [87--100\%]; bar $\ge$ 50\%): 25 equal within $10^{-6}$, 1 above by $1.2\times10^{-6}$ & 26/26 \\
Anchor parent (the model's own anchor) & 28 & 7/28 (25\%, [13--43\%]); per cell 4/12, 1/5 and 2/11 at $N = 13$, 21 and 31 & P-MU1 predicted $\ge$ 50\%, not met & 17/28 above the anchor, 4/28 below & 22/28 \\
Off-family parent (diagonal chain, sum $\approx$ 0.6) & 34 & 20/34 abandon parent for $k^*$-grid (59\%, [42--74\%]); per cell 1/7 [3--51\%], 7/13 [29--77\%] and 12/14 [60--96\%] at $N = 13$, 21 and 31 & P-MU4: $k^*$-grid reasserts, scored per cell & --- & --- \\
\bottomrule
\end{tabular}

{\footnotesize Bracketed ranges are 95\% Wilson intervals on the rate they follow.}
\end{small}
\end{center}


\textbf{Clearance of the family argmax (post hoc, disclosed).} Arm MU's rows were scored
against the anchor, so the clearance question was never asked at registration. Recomputing
the argmax from the closed form gives $1.776142375$ at $N = 13$. At $N = 21$ it is
$2.258883476$. At $N = 31$ it is $2.748528137$. The 135 invocations give 88 valid outputs at
the primary tolerance. At $10^{-9}$ just 63 of the 135 are valid. Of the 88, 27 return a
sum nominally above the argmax: 9 within the $5\times10^{-8}$ granularity of a recorded sum,
18 beyond it, and 61 at or below. Seventeen of the eighteen are argmax-parent outputs
returning their parent's score within $1.2\times10^{-6}$. The eighteenth is the anchor-parent output at
$N = 21$ that moved onto the argmax.

The argmax parent's seven-decimal coordinates sum to $3.25\times10^{-7}$ above the
closed form at $N = 13$ and $2.63\times10^{-7}$ above it at $N = 31$. Thirteen of the
eighteen return that parent sum unchanged; their excess is the parent's rounding. The other
five sit up to $1.53\times10^{-6}$ above the argmax. All eighteen fail the $10^{-9}$ test on
overlapping circles, not on containment. Their largest single overlaps run from
$3.8\times10^{-9}$ to $4.3\times10^{-7}$, median $4.4\times10^{-8}$, each smaller than its
row's excess. Rounded or perturbed coordinates carry the excess; the $10^{-6}$ tolerance
admits the overlaps they create (per-row table and the granularity split of the 88,
Supplement~S2). Of the 63 valid at $10^{-9}$, four sit above the argmax by at most
$2.5\times10^{-8}$, inside the recorded-sum granularity, and none clears.

\section{Five loop generations: arm L}\label{sec:arm-l}

Arm L iterates arm MU's conditioning step for five generations, feeding each scored proposal
back as the next parent. Iteration keeps the proposer off the anchor: 11 of 49 conditioned outputs sit on it,
against 11 of 17 at generation 0. Three of the four lineages end exactly on the family argmax
and none goes above it.

\textbf{Design.} Both prompts were hashed, and the stepper and scorer were committed
alongside. Generation 0 is the bare prompt, sampled
fresh inside this arm. Its prompt is byte-identical to the companion paper's unconditioned
arms, served on the same channel and wrapper. Generations 1 to 5 use arm MU's
registered mutation prompt verbatim, instantiated with the parent packing and its score. Two
discriminating cells, $N = 13$ and $N = 31$, are crossed with two archive regimes. GREEDY
retains only the single best valid proposal. DIVERSE retains up to five distinct-valued ones.
That gives four lineages of six rounds at population 5, which is 120 invocations. Two of the
120 were rejected by the runtime before reaching the model, and are counted as invalid rather
than dropped. The parent for population slot $i$ is
$\text{archive}[i \bmod |\text{archive}|]$, so the loop carries no sampling randomness of its
own. All four lineages are evaluable under the registered floor.

This is a minimal analogue. It has no islands, no crossover and no program representation.
Its only evaluator signal is the parent's own score.

\textbf{Registered predictions.} L1 predicted that the modal valid output at generation 0
equals the anchor in every lineage. L2 predicted that dissolution survives
iteration, so that the pooled on-prediction rate over generations 1 to 5 falls below the
generation-0 rate. L3 predicted that the pooled conditioned on-prediction rate is higher under
GREEDY than under DIVERSE, on the reading that loop diversity is scaffold-produced. L4 was
registered open with no direction predicted: report best-of-run against the family argmax.
The falsifier fires if the pooled conditioned rate meets or exceeds the generation-0 rate in
both cells.

\begin{table}[!ht]
\caption{Arm L by lineage. Valid outputs over invocations, and on-prediction counts, at
generation 0 and across the five conditioned generations, with each lineage's best scoring
valid output, the generation at which that best first appeared, and its cell's family
argmax. Pooled: 11 of 17 on-prediction at generation 0, 65\% (95\% Wilson 41--83\%), against
11 of 49 conditioned, 22\% (13--36\%).}
\label{tab:arm-l}
\begin{center}
\footnotesize
\setlength{\tabcolsep}{4pt}
\begin{tabular}{@{}lccccccc@{}}
\toprule
Lineage
& \begin{tabular}[b]{@{}c@{}}Gen-0 valid\\/ invocations\end{tabular}
& \begin{tabular}[b]{@{}c@{}}Gen-0\\on-prediction\end{tabular}
& \begin{tabular}[b]{@{}c@{}}Conditioned valid\\/ invocations\end{tabular}
& \begin{tabular}[b]{@{}c@{}}Conditioned\\on-prediction\end{tabular}
& \begin{tabular}[b]{@{}c@{}}Best\\of run\end{tabular}
& \begin{tabular}[b]{@{}c@{}}Best first\\reached (gen)\end{tabular}
& \begin{tabular}[b]{@{}c@{}}Family\\argmax\end{tabular} \\
\midrule
13-GREEDY & 5/5 & 2 & 4/25 & 0 & $1.7761424$ & 0 & $1.7761424$ \\
13-DIVERSE & 4/5 & 2 & 16/25 & 5 & $1.7761424$ & 5 & $1.7761424$ \\
31-GREEDY & 4/5 & 3 & 18/25 & 3 & $2.6104167$ & 5 & $2.7485281$ \\
31-DIVERSE & 4/5 & 4 & 11/25 & 3 & $2.7485281$ & 2 & $2.7485281$ \\
\midrule
Pooled & 17/20 & 11 & 49/100 & 11 & & & \\
\bottomrule
\end{tabular}
\end{center}
\end{table}

\textbf{Validity, and what every rate in this arm conditions on.} Validity drops under
iteration. The pooled row of Table~\ref{tab:arm-l} gives the two totals. The weakest lineage,
13-GREEDY, carries the loss. Conditioned rounds lose validity as well as accuracy, so every
rate here counts only survivors. The 11 of 49 on-prediction rate is one such rate. Over all
invocations it is 11 of 100. That is 6--19\%. At generation 0 it was 11 of 20. That is
34--74\%. The registration set no validity prediction, so the loss is a post hoc
observation.

\textbf{L2 confirmed, with unequal support.} Sum the two lineages at each cell of
Table~\ref{tab:arm-l}. At $N = 31$ the generation-0 on-prediction rate is 7 of 8. Conditioned,
it falls to 6 of 29. That is 21\% (10--38\%). At
$N = 13$ it starts at 4 of 9. Conditioned, it is 5 of 20. That is 25\% (11--47\%). Both drops
are point estimates. The falsifier did not fire. Once a lineage's parent sits
off the anchor, later rounds condition on a parent at the family argmax, arm MU's second
condition, so L2 is not independent of MU. Of the 100 conditioned rows, 34 carry a parent at
the family argmax: 25 in 13-GREEDY and 9 in 31-DIVERSE, none in the other two lineages.
13-GREEDY returned 4 of 25 valid under that
condition where arm MU's argmax-parent rows returned 26 of 45, a gap this paper does not
explain; 31-DIVERSE returned 7 of 9.

\textbf{L1 was not met.} It fails on one lineage. The 13-DIVERSE seed round in
Table~\ref{tab:arm-l} splits evenly, and the registered tie rule counts the tie against the
prediction. The other three lineages pass. The rule set four per-lineage bars rather than one
pooled bar per cell. Each bar ran over the four or five valid proposals that lineage's seed
round returned.

\textbf{L3 was not met.} Per cell, GREEDY and DIVERSE give 0 of 4 against 5 of 16 at
$N = 13$ and 3 of 18 against 3 of 11 at $N = 31$. Pool each archive regime over both cells.
The conditioned
on-prediction rate is 3 of 22 under GREEDY. That is 14\% (5--33\%). Under DIVERSE it is 8 of
27. That is 30\% (16--48\%). The difference runs opposite to the
registered direction and is not separable from chance. A post hoc Fisher exact test on the
pooled $2 \times 2$ table gives $p = 0.303$. The odds ratio is 0.375. The comparison is also
confounded: lineages differ in parent identity as well as archive width. 13-GREEDY
also contributes the fewest conditioned valid proposals of any lineage. The outcome is a
disconfirmation of L3, not evidence for its converse.

\textbf{L4: three of the four lineages reach the family argmax exactly, and none clears it.}
The best-of-run and family argmax columns of Table~\ref{tab:arm-l} carry the comparison. Both
lineages at $N = 13$ and 31-DIVERSE land exactly on their cell's family argmax, and 31-GREEDY
stops short of it. The best of 13-GREEDY is a generation-0 output, one family-argmax packing in its unconditioned seed round that
the loop never bettered. The other three lineages reach their best under conditioning. Of the 49 valid
conditioned outputs the table counts, 47 are valid at $10^{-9}$, and none of those exceeds
the best value the recipe family can express by more than $10^{-6}$. The largest excess among
the 49 is $1.3\times10^{-15}$, float noise, on five 31-DIVERSE rows. One post-sampling scorer
correction, affecting only the
$10^{-9}$ test and no verdict, is disclosed in Supplement~S2.

\section{Serving path and reasoning budget: arms P and P-D}\label{sec:arm-p}

Arms P and P-D change how the proposer is served, not what it is conditioned on. On the bare
pinned path validity collapses, 0 of 105 at the square cells, so no mode prediction can be
scored. With extended thinking on, validity returns, 12 of 14, and the anchor is again the
plurality value. Per-cell tables sit in Supplement~S5.

\textbf{Arm P: the registered prompt on a bare, pinned serving path.} Arm P sends the
registered prompts byte-identical through a path that exposes and pins temperature. It
runs at $n = 15$ per cell across seven square cells, five held-out cells and three code cells. That
is 225 invocations. The serving path refused 6, none at a square cell, each counted invalid
where it was issued. Two competing mode predictions say the anchor $T(k^*, N)$ is modal in a
majority of the seven square cells, or in few of them. Two competing code-channel predictions
say no valid program clears the family argmax, or at least one does. The falsifier is on the
discriminating-cell rate. Supplement~S1 states all five in the registration's own words.
The registration left 4 of the 7 square cells that neither mode prediction adjudicates, a
design gap. An outcome there confirms neither.

\textbf{Validity collapsed.} The seven square cells carry both mode predictions. There, 0 of
105 outputs are valid at the primary tolerance, so none at $10^{-9}$. All 105 returned well-formed lists of the requested
count, with no prose, fences or reasoning tokens. Every failure but three is an
overlap. Supplement~S5 names the three. The median maximum pairwise overlap at $N = 13$ is
0.10. That is fifty times the value window. At the held-out cells 4 of 75 are valid, all four
at $N = 50$. They sum to 2.5, a full unit below the family argmax there, so none clears.

\textbf{Against the registered floor.} The arm's own floor is five valid outputs per cell, so
the direct-emission predictions are UNSCOREABLE at every cell. The falsifier cannot fire either. Its
denominator, the valid outputs at the four discriminating square cells, is
empty. Empty denominators receive no imputed rate. On the
code cells 12 of 45 programs are valid, and only $N = 21$ reaches the floor. None clears the
family, so the no-clearance prediction holds.

\textbf{Two post hoc diagnostics bound the cause.} Neither was registered. The first is a
temperature sweep. At 0.0, 0.5 and 1.0 the $N = 13$ prompt gives 0 of 6 valid per setting. The
second is a same-day control at the same cell through the agent runtime, tabled beside the two
arm P-D conditions in Supplement~S5.

\textbf{Arm P-D: which component of the instrument.} Arm P-D was registered after an outside
review asked for a diagnostic beyond temperature. It sends the same $N = 13$ prompt, to the same alias, on the same pinned path and
decoding, under two one-line manipulations, at $n = 15$ each. The \emph{thinking-on}
condition turns the vendor's extended thinking on at the routing layer's default effort and
sends no system prompt. The \emph{thinking-off} condition leaves thinking off and adds one
system line asking for care. The reasoning-budget
hypothesis predicts thinking-on valid in at least 8 of 15 against at most 3 for thinking-off.
The context hypothesis exchanges the two thresholds. A third registered outcome covers
anything else. The arm closed at 27 of 30 rows by a registered amendment, after the serving
path refused the last three for credit. Those rows could not have changed the verdict,
because thinking-off could reach at most 2 valid while thinking-on already held 12.
Thinking-on is valid in 12 of 14. The 95\% Wilson interval on that rate is 60--96\%.
Thinking-off is valid in 0 of 13, every failure an overlap, exactly arm P's collapse. So the
reasoning-budget hypothesis holds, and the anchor is the plurality value. $T(4,13) = 1.625$
accounts for 5 of the 12 valid thinking-on outputs. The registered secondary prediction, that
this value is modal in whichever condition clears the validity threshold, is therefore met.
The two conditions differ in thinking and in the system line, but arm P sent this prompt with
thinking off and no system line and collapsed, so thinking-on changes one feature against that
baseline. Only $N = 13$ was run, so the component is named on one cell, and no interval is
computed there.

\section{Other vendors: arms GM, GM2, GM3 and V}\label{sec:gm3-v}

Arms GM, GM2, GM3 and V change the vendor. GM3, the only chain arm with a verdict, meets the
registered pooled bar, 46 of 80. In a post hoc split it misses at the discriminating cells,
11 of 43, where the modal output is the family argmax construction. In arm V ten of thirteen aliases fall below the floor, and
the three that clear it split, two anchoring at point estimate and one not. Arms GM, GM2 and
GM3 run against a second vendor's first-party API and vary the model family, with prompts
pinned, an output budget fixed and a provider logged per row. Arm V varies vendor and serving
class at once.

\textbf{Arms GM and GM2 returned no scoreable data, and both are reported.} One chain
registration fixes the definitions and predictions for all three GM arms. Each arm addresses
seven $N$ cells by twenty samples, over a floor of three valid samples per cell. The primary
prediction requires a mode match in at least 5 scoreable cells. The second requires a pooled
on-prediction rate of at least 30\%. The third requires the two-radius signature at order
$k^*$ in at least half of on-prediction samples. That signature is a grid radius $1/(2k^*)$
carrying a filler radius $(\sqrt{2} - 1)/(2k^*)$. The falsifier fires on a mode-match failure
in at least 4 scoreable cells. Arm GM addressed a gemini-flash-lite alias, part of it through
a routed path (Supplement~S2). It returned 20 valid of 140, UNDERPOWERED under both
accountings, and its falsifier fires under neither. Arm GM2 addressed the gemma family at a 4,096-token output budget, and returned 0 of
140 parseable. The cause is a budget confound, not a model result. The vendor API returns
deliberation on a separate channel, and the budget was spent there before any answer text. Arm
GM3 re-ran the design at 16,384 tokens. Of its 140 rows, 17 did not parse. Of the 123
that did, 80 are valid packings at $10^{-6}$.

\textbf{Arm GM3: the pooled bar is met and the discriminating cells are not.} Arm GM3 runs the
unconditioned-call protocol against gemma-4-26b through the vendor's first-party API. The
first-party key reached its quota, so 41 of the 140
rows were sampled through a routed path under an amendment committed before resumption, and
both accountings are reported.

The pooled prediction held. Of the 80 valid packings, 46 land within $2\times10^{-3}$ of the
value the closed form predicts. That clears the registered pooled bar. The 41 routed rows
carry 2 of the 80 valid packings. None of the 46 on-prediction are among them. First-party
rows alone
therefore give 46 of 78, and no verdict changes. A pooled rate cannot separate family search
from anchoring, because the two agree off the discriminating cells. The two cells where the anchor equals the family argmax
hold 35 of the 46 on-prediction packings. On the three scoreable discriminating cells, a
split read post hoc, the rate is 11 of 43, below the bar the pooled metric passes.

The cell-level bar was not met. The anchor is modal in 2 of 5 scoreable
cells, and Figure~\ref{fig:gm3anchoring} plots the per-cell modal sums against the
anchor. The falsifier did not fire either: three failures, short of its four. Bar
and falsifier together leave one to three failures unadjudicated, a registration design gap,
and the observed three fall in it. The two-radius signature also came in under its bar of half. Only 20 of 46 on-prediction
samples carry the exact decomposition. That bar cannot be met at a truncate cell, whose anchor carries one radius
and no filler. The 15 on-prediction samples at $N = 35$ are counted against a signature they
could not show.

\begin{figure}[t]
\centering
\begin{minipage}[c]{0.44\linewidth}
\centering
\includegraphics[width=\linewidth]{fig_gm3_anchoring.png}
\end{minipage}\hfill
\begin{minipage}[c]{0.53\linewidth}
\caption{\emph{Arm GM3 modal sums against the anchor.} Filled green circles
match the anchor (the dashed line is $V(k^*, m)$ at extend cells and $T(k^*, N)$ at
truncate cells), with frequencies annotated. Open red triangles miss, and every miss lands
above the anchor at the family argmax \emph{value}. The post hoc structural check in the
text, not this figure, establishes that these are also the family argmax \emph{construction}.
Grey crosses on the anchor line are the two UNSCOREABLE cells, with their valid counts.
Regenerates via \texttt{fig\_gm3\_anchoring.py}
from \texttt{arm\_gm3\_report.json}.}
\label{fig:gm3anchoring}
\end{minipage}
\end{figure}

\textbf{The misses are the family argmax construction (post hoc, disclosed).} The modal
output at every discriminating cell is the family argmax itself. A post hoc structural diagnostic finds all 27 of the 27 packings at the family
argmax value carrying its two-radius decomposition, at the $10^{-3}$ radius tolerance the
diagnostic uses. That
decomposition is the $(k^*{-}1)$ grid radius with the $(k^*{-}1)$ filler radius. The
registered signature test could not see this, because it tests radii against the predicted
order $k^*$. On the discriminating cells 27 of 43 valid packings are the family argmax
construction. That is 63\%, on a 48--76\% interval. Validity collapses at the two largest
cells, $N = 37$ and 43, which bounds where this model returns evaluable packings; both are
UNSCOREABLE under the registered rule. Clearance was re-read post hoc,
as for arm MU. Of the 80 valid packings, 79 are valid at $10^{-9}$. None clears the family
argmax under the rule of \S\ref{sec:method}. Twenty-nine of the 79 sit above it by at most
$5.3\times10^{-10}$, the rounding of a printed coordinate. The rest sit at or below
(Supplement~S4).

\textbf{Arm V: thirteen free-tier aliases.} Arm V sends the GM chain's prompt files
byte-identical at five cells, $N = 13$, 17, 31, 35 and 37. It draws five samples per model and
cell. The targets are free-tier aliases across seven attributed vendors. Three amendments,
each committed before the sampling it
governs, expanded the original three-alias design to thirteen aliases across two serving
paths. The last raised the output budget after hidden-reasoning truncation emptied three
models' replies.

Execution retried each invocation slot until content arrived, per the registered
transport-repair rule. The per-model floor of \S\ref{sec:method} is evaluated on registered
slots, so a slot with no row counts against it. Ten of thirteen aliases finished BELOW-FLOOR,
including all three originally registered aliases. One of them made 108 attempts across its 25
registered slots, none of them with content. The other three cleared the floor, and their
registered verdicts split.

V1 predicted at least 50\% of a model's valid outputs on-prediction, and V2 at most 10\% at
the family argmax. V2 is scored only at the discriminating cells, $N = 13$ and 31. In this
arm, \emph{transfer} denotes only the registered pooled-bar verdict. An alias must clear the
per-model floor and meet the V1 bar on its valid outputs pooled over five cells. Every transfer
verdict here rests on amendment-added aliases. All three were registered by amendment 2
(\texttt{arm\_v\_\bk{}preregistration\_\bk{}amendment2.md}, \texttt{d579797}).

\begin{center}
\footnotesize
\setlength{\tabcolsep}{3.5pt}
\begin{tabular}{@{}llcccc@{}}
\toprule
alias & vendor & valid & V1 on-prediction & V1 bar ($\ge$50\%) & \begin{tabular}[b]{@{}c@{}}V2 family argmax\\valid at discriminating cells\\(bar: $\le$10\%)\end{tabular} \\
\midrule
\texttt{north-mini-code} & Cohere & 12 & 8/12 [39--86\%] & \textbf{MET} (point est.) & 0/4 \textbf{HOLDS} \\
\texttt{gpt-oss-20b} & OpenAI & 18 & 11/18 [39--80\%] & \textbf{MET} (point est.) & 0/9 \textbf{HOLDS} \\
\texttt{gemma-4-31b-it} & Google & 15 & 0/15 = 0\% [0--20\%] & \textbf{NOT MET} & 0/3 \textbf{HOLDS} \\
\bottomrule
\end{tabular}

{\footnotesize Bracketed ranges are 95\% Wilson intervals on the rate they follow. V2 is
scored only at the discriminating cells. The registration labels a met V1 bar TRANSFERS and a
missed one DOES-NOT-TRANSFER; this table reports the bar outcome, in the sense
\S\ref{sec:gm3-v} defines. Both passing intervals span the 50\% bar.}
\end{center}


\textbf{Both passes are weaker than the percentages suggest.} The
bracketed Wilson 95\% intervals include the bar, so each verdict is a point-estimate pass.
And the pooled rate spans cells that cannot
discriminate anchoring from family search. Decomposing onto the discriminating cells gives 3
of 4 for the first alias. It gives 5 of 9 for the second. The arm's strongest reading, that these aliases behave
as the original family does, does not survive the template-shape null. That null is a proposer
picking uniformly at random among the admissible grid orders in $k^*{-}1$, $k^*$ and $k^*{+}1$,
one branch value each. It lands on-prediction one time in three at every cell but $N = 17$,
where one order overfills its interstices and the rate is one in two. Supplement~S5 derives it in
this paper's own terms and gives the exact tails, and this paper applies it post hoc
\citep{companion}. The pooled rates clear that null only because most of their cells cannot
discriminate. At the cells that can, neither pass separates from it. Arm GM3's
post hoc 11 of 43 at its discriminating cells does not separate from it either. V2 is a point-estimate
pass on 16 valid outputs at discriminating cells: 0 land on the family argmax, on per-alias
counts of 3, 4 and 9 whose Wilson intervals all include the 10\% bar. The same post hoc re-read over the three evaluable aliases
finds 45 valid at $10^{-6}$ and 39 at $10^{-9}$. None clears the argmax; eight of the 39
sit above it by $8.9\times10^{-16}$, float noise, and the rest sit at or below. The closest at
a discriminating cell
falls short of it by $2.5\times10^{-3}$ at $N = 31$ (Supplement~S4).

\textbf{Anchoring is not confined to the original family, and not universal.} Two further
vendors' models land the registered construction values at point estimates above the bar.
Through the chain's first-party API, gemma-4-26b emits the family argmax construction at
the discriminating cells. Through arm V's free-tier serving path, gemma-4-31b lands neither the
anchor nor the family argmax. Its valid packings are de-novo, unanchored and weak. Their per-cell bests run
0.88--1.85. The anchors at those cells run 1.63--3.03. Model variant and serving path
change together across this pair, so the boundary is located but not attributed.

\section{Related work}\label{sec:related}

\textbf{The loop the proposer is characterized for.} Language Model Crossover
\citep{meyerson2024languagemodelcrossovervariation}, ELM
\citep{lehman2022evolutionlargemodels}, EvoPrompt \citep{guo2024evoprompt}, EoH
\citep{liu2024eoh} and LLaMEA \citep{vanstein2025llamealargelanguagemodel} made the LLM call
an evolutionary operator over prompts and programs. FunSearch
\citep{romeraparedes2024funsearch} turned the pattern into a discovery claim. AlphaEvolve
\citep{novikov2025alphaevolve} and its successors inherit that claim, OpenEvolve
\citep{sharma2025openevolve} among them
\citep{lange2025shinkaevolveopenendedsampleefficientprogram,su2026helixevolutionaryreinforcementlearning,khrulkov2025gigaevoopensourceoptimization,cemri2026adaevolveadaptivellmdriven,wang2025thetaevolvetesttimelearningopen}.
Every one conditions its proposer on a parent and iterates it, two of the conditions this
paper varies. Simple baselines \citep{gideoni2026simplebaselinescompetitivecode} and
classical solvers \citep{berthold2026outoftheboxglobaloptimizationpacking} recover much of
the reported advantage.

\textbf{Template convergence under iteration.} ``Mutation Without Variation''
\citep{gurkan2026mutationvariationconvergencedynamics} finds iterated LLM program mutation
collapsing onto previously seen structural templates. In 87\% of chains, over 93\% of
mutations revisit a seen form. It runs mutation loops in program space with no closed form
and no preregistration, and is the closest published neighbour to arm L. Its result is about
structural revisiting. Arm L's second prediction is about a value the revisited structures
do not cross.

\textbf{Anchoring, and what kind of object anchors here.} Anchoring is studied as a general
bias over initial information, and as numeric primes dragging point estimates
\citep{huang2026understandinganchoringeffectllm,lou2024anchoringbiaslargelanguage,malberg2025comprehensiveevaluationcognitivebiases}.
The modality here differs: a construction template has a computable value, where a scalar
anchor has only a measurable pull. One level up, LLM outputs concentrate where human ones do
not \citep{chen2026measuringgaphumanllm,jiang2025artificialhivemindopenendedhomogeneity}.
Set-level and diversity-aware selection
\citep{li2026dictionariesdarwinsetlevelselection,wang2026diversescientifichypothesissearch,yang2026accuracyevaluatingstrategydiversity}
and the bin-packing critiques
\citep{herrmann2026indepthstudyllmcontributions,sim2025hypebenchmarkingllmevolvedheuristics}
converge on the same question about the proposer's contribution.

\textbf{Two results that cut against this paper's framing.} Evolutionary search matters
empirically \citep{zhang2024understandingimportanceevolutionarysearch}, so one unconditioned
call is a poor stand-in for a loop. Strong LLM optimizers act as local refiners
\citep{zhang2026makesllmgoodoptimizer}, as a parent-conditioned call does.

\textbf{Instrument effects and reasoning budget.} Larger and more instructable models become
less reliable \citep{zhou2024reliability}. A reasoning model can report a higher
hallucination rate alongside higher accuracy \citep{openai2025o3o4systemcard}, and reasoning
collapses past a complexity threshold
\citep{shojaee2025illusionthinkingunderstandingstrengths}. Geometric construction is a regime
where nominal capability and executed correctness diverge
\citep{kim2026geobuildbenchbenchmarkinteractiveexecutable,balunovic2025mathconstructchallengingllmreasoning}.

\textbf{Preregistration lineage.} Recipes, outcome-blinded predictions and agent protocols
are already preregistered
\citep{thomas2026mitigatingllmbasedphackingpreregistering,williams2026predictingllmsafetyrelease,vaccaro2026preregistrationexperimentsaiagents,zhang2026testingfrontierlargelanguage}.
HindsightBench \citep{jia2026hindsightbenchblackboxbehavioralaudit} releases frozen
preregistrations of directional aggregate hypotheses. This paper's arms lock a numeric
threshold and a falsifier before each setting was changed.

\section{Limitations, reproducibility, deviations}\label{sec:limits}

\textbf{Scale.} Arm MU takes one conditioning step, and arm L runs five generations at
population 5 with one mutation operator. A larger loop, with islands and crossover, could
sustain variation this one cannot.

\textbf{No instruction-free control.} No call without the improve instruction was run. So
returning the argmax parent unchanged,
26 of 26, cannot be split into recognising an optimum and ignoring the instruction. Only its
parent-dependence is established: the same instruction moves 17 of 28 off the model's own
anchor.

\textbf{Vendor and tier scope.} The proposer under test is one vendor's weak-tier alias,
except where another is named. The cross-vendor arms move the pooled-level statement, not the
cell-level mode structure.

\textbf{Prompt robustness.} Every arm conditions on a single hash-locked prompt. That buys
reproducibility, not robustness to rewording. The companion paper reworded two discriminating
cells and found the same modal construction \citep{companion}.

\textbf{Untested alternatives.} Typicality bias in preference-tuned models
\citep{zhang2026verbalizedsamplingmitigatemode} predicts concentration on the most typical
construction without reference to geometry, and nothing here excludes it. Contamination
probes \citep{carlini2021extracting,golchin2024timetravel}, a format-sensitivity sweep
\citep{sclar2024formatspread} and alignment-induced diversity loss \citep{kirk2024rlhf} remain
untested.

\textbf{Reproducibility.} Every arm's prompts were hashed, and its predictions and falsifiers
fixed, before sampling. Exceptions are in the supplement's deviation table, and no arm was
selected on its result. The anonymized artifact, shipped as supplementary material with entry
point \texttt{artifact/BUNDLE\_README.md}, holds the ledgers, scripts,
registrations, amendments and frozen reports. It regenerates every reported count by a
documented command, and a replay gate (\texttt{replay\_gate.py}) runs each documented
script from the packaged bundle. Arms P and P-D can be re-sampled on their pinned path. Arms MU and L can be re-scored
but not re-drawn, because their rows depend on a dispatch wrapper and inherited context the
digests do not lock.

\textbf{Where the forking paths are.} Five choices were made with data already in hand, and
each is named here. (a) Arm V's evaluable set exists because of amendment 2
(\texttt{d579797}), which expanded the alias list after the three registered aliases finished
below the floor. (b) Arm P-D was designed after arm P collapsed, and registered at
\texttt{5554cc0}. (c) The clearance readings for arms MU, GM3 and V, GM3's
discriminating-cell split and arm L's pooled Fisher test are post hoc, registered for none. (d) Both arm P diagnostics, the temperature sweep and the same-day runtime control,
were chosen after that collapse. (e) Arm P's scoring code was committed at \texttt{596b113},
after its registration at \texttt{aa473ad} and after sampling. None is corrected for
multiplicity. The protection is that each is listed.

\textbf{Deviations, in summary.} Arms GM and GM3 each carry
one registered serving-path deviation, with both accountings reported for each, and GM3 also
carries a disclosed duplicate-row repair. Arm V has three pre-sampling amendments. Arm L has
one post-sampling scorer correction, affecting only the $10^{-9}$ clearance test. Arm P-D
closed at 27 of 30 rows by an amendment that could not change its verdict.

\textbf{Use of AI systems.} Claude models drafted prose and assisted implementation. The
authors designed the study, specified the instruments and their decision rules, and approved
each preregistration before sampling. Each such commit is a git ancestor of the sampling it
governs.

\section{Conclusion}\label{sec:conclusion}

Eight registered arms changed four things about the setting, for one task and one weak-tier
proposer. On the agent-runtime instrument, one conditioning step under the improve instruction
left 0 of 26 outputs on the anchor given a better parent. An off-family parent brought the
$k^*$-grid back in 20 of 34. Five loop generations there did not restore the anchor. A pinned serving path collapsed validity until extended
thinking was on. At a second vendor the pooled bar held while the discriminating cells did
not. A zero-shot characterization therefore
does not establish how a proposer behaves inside a discovery loop without a measurement at
that loop's own conditions. The boundary is two-sided: the anchor does not carry into a
conditioned call, the family argmax does, as a ceiling no arm clears under the rule of
\S\ref{sec:method}. Check coordinates as well as score:
a familiar score can come from a different arrangement.


% Citations use natbib \citet/\citep per TMLR's instructions.
\bibliography{references}
\bibliographystyle{tmlr}

\end{document}
